\section{Deployment}

\subsection{Deployment Model}
AgriApp deployment is structured as a repeatable operations workflow, not a one-time installation. The Raspberry Pi is the edge runtime target and hosts backend APIs, camera orchestration, and local data storage. Development happens on a workstation, then changes are synchronized incrementally to the device.

This model is intentionally lightweight: rapid iteration is required for hardware/software co-tuning, and full image reflashing would be too costly for daily work.

\subsection{Primary Tooling: \texttt{deploy\_rpi.sh}}
The project standardizes deployment operations through \texttt{deploy\_rpi.sh}. In current practice, this script is used for four classes of tasks:
\begin{itemize}
  \item synchronization and remote access (\texttt{--sync}, \texttt{--mount}, \texttt{--ssh}),
  \item backend runtime reload (\texttt{--reload-server}),
  \item ESP build/flash support (\texttt{--esp-build}, \texttt{--esp-flash}),
  \item Android build/install/run pipeline (\texttt{--android-build}, \texttt{--android-install}, \texttt{--android-run}, \texttt{--android-cir}).
\end{itemize}

This unifies day-to-day engineering actions behind one operator-facing entry point.

\subsection{Raspberry Pi Service Deployment}
Backend serving is managed as a user-level systemd service (\texttt{agriapp-server.service}). Installation is performed by \texttt{scripts/install\_systemd\_user\_services.sh}, and lifecycle is controlled through \texttt{systemctl --user}.

Using systemd for the server gives deterministic startup behavior after reboot and clear operational observability (status and logs). It also allows remote restart through dedicated API/system actions.

\subsection{Capture Runtime Deployment}
Current deployment practice separates backend service supervision from capture scheduling logic:
\begin{itemize}
  \item server process: supervised by systemd,
  \item unattended capture loop: controlled dynamically via API endpoints (start/stop/status).
\end{itemize}

This split was chosen to provide direct operational control from tablet/web interfaces without editing cron/systemd timers during field sessions.

\subsection{Environment and Dependencies}
Runtime wrappers (notably \texttt{run\_server.sh} and capture wrappers) support pyenv/venv-based Python environment loading. The deployment checklist must ensure that required packages are present on the Raspberry Pi, including image handling and optional sensor libraries.

For sensor-enabled metadata enrichment, deployment validation should include:
\begin{itemize}
  \item GPS visibility on UART serial path (\texttt{/dev/serial0}),
  \item BME280 visibility on I2C bus,
  \item successful metadata sidecar generation after one-shot capture.
\end{itemize}

\subsection{GPS UART Wiring Reference}
For compact GPIO wiring, GPS can be connected through the hardware UART:
\begin{itemize}
  \item \textbf{TX} (GPS) $\rightarrow$ \textbf{RPi RXD0 / GPIO15} (pin 10),
  \item \textbf{RX} (GPS) $\rightarrow$ \textbf{RPi TXD0 / GPIO14} (pin 8),
  \item \textbf{GND} (GPS) $\rightarrow$ \textbf{RPi GND} (pin 14),
  \item \textbf{VCC} (GPS) $\rightarrow$ \textbf{RPi 3.3V} (pin 17).
\end{itemize}

During deployment, UART console must be disabled on \texttt{serial0} and hardware UART enabled (\texttt{enable\_uart=1}) to avoid conflicts with Linux boot/login console.

\subsection{BME280 Wiring Reference}
In I2C mode, the minimal wiring uses four lines:
\begin{itemize}
  \item \textbf{VCC} $\rightarrow$ \textbf{RPi 3.3V} (pin 1),
  \item \textbf{GND} $\rightarrow$ \textbf{RPi GND} (pin 6),
  \item \textbf{SDA} $\rightarrow$ \textbf{GPIO2 / SDA1} (pin 3),
  \item \textbf{SCL} $\rightarrow$ \textbf{GPIO3 / SCL1} (pin 5).
\end{itemize}

On six-pin boards, \texttt{CSB} and \texttt{SDO/SDD} are optional for basic I2C operation and can be used only when address forcing is needed.

\subsection{Android Deployment Cycle}
Android deployment is integrated into the same repository workflow. The standard cycle is build-install-run on USB-debug device via \texttt{--android-cir}. Validation after install must include discovery, one-shot capture actions, gallery pagination, and system/network controls.

\subsection{Network Portability}
The edge node is expected to move across multiple Wi-Fi environments (e.g., home and field networks). Deployment therefore includes maintaining multiple known SSIDs with proper priority, plus AP rescue fallback for infrastructure outages.

Operationally, AP mode is not a special-case debug path: it is part of the primary resilience strategy and should be tested regularly.

\subsection{Recommended Deployment Runbook}
A practical release cycle is:
\begin{enumerate}
  \item local validation (backend syntax, Android build, LaTeX if docs changed),
  \item \texttt{deploy\_rpi.sh --sync},
  \item \texttt{deploy\_rpi.sh --reload-server},
  \item API smoke checks (health, status, gallery, one-shot),
  \item optional tablet functional checks.
\end{enumerate}

This sequence has proven robust in both office and field transitions.

\subsection{Scenario Runbooks}
To reduce ambiguity during field operation, deployment is documented as scenario-specific runbooks.

\paragraph{Scenario A: Studio code update + RPi rollout}
\begin{enumerate}
  \item local smoke checks (\texttt{py\_compile}, critical scripts),
  \item \texttt{./deploy\_rpi.sh --sync},
  \item \texttt{./deploy\_rpi.sh --reload-server},
  \item verify \texttt{/api/v1/health} and one-shot capture.
\end{enumerate}

\paragraph{Scenario B: RPi reboot after power cycle}
\begin{enumerate}
  \item verify user service status (\texttt{systemctl --user status agriapp-server.service}),
  \item verify network reachability (expected Wi-Fi profile or AP fallback),
  \item run one-shot capture and gallery check.
\end{enumerate}

\paragraph{Scenario C: Sensor troubleshooting}
\begin{enumerate}
  \item stop backend service to avoid serial contention,
  \item run \texttt{test/test-gps.py} and \texttt{test/test-bme.py},
  \item restart backend service,
  \item validate metadata sidecar generation on a fresh shot.
\end{enumerate}

\paragraph{Scenario D: SegGPT comparative experiments}
\begin{enumerate}
  \item align ML stack versions (Studio and RPi),
  \item verify same image set and prompt mask availability on both nodes,
  \item run \texttt{scripts/seggpt\_compare\_benchmark.py},
  \item archive \texttt{summary.csv} and JSON traces for paper artifacts.
\end{enumerate}

\subsection{Version-Parity Baseline for Reproducibility}
For current comparative segmentation campaigns, the operational baseline requires aligned versions across Studio and RPi:
\begin{itemize}
  \item \texttt{torch 2.10.0},
  \item \texttt{transformers 5.5.0},
  \item \texttt{torchvision 0.25.0} (Studio where applicable).
\end{itemize}

Without parity, box/mask outputs can diverge even when using identical images and prompt masks, which invalidates strict reproducibility claims.

\subsection{Unified Version Matrix (Current Baseline)}
\begin{longtable}{p{0.26\linewidth} p{0.22\linewidth} p{0.22\linewidth} p{0.22\linewidth}}
\toprule
Component & Studio Node & RPi Node & Note \\
\midrule
Python runtime & \texttt{/home/fra/pyvenv} & \texttt{/home/fra/pyenv} & isolated venv per node \\
Torch & \texttt{2.10.0} & \texttt{2.10.0} & parity required for SegGPT comparisons \\
Transformers & \texttt{5.5.0} & \texttt{5.5.0} & parity required for mask/box equivalence \\
Torchvision & \texttt{0.25.0} & optional / not required in baseline path & mismatch can break model imports on Studio \\
SegGPT checkpoint & \texttt{BAAI/seggpt-vit-large} & \texttt{BAAI/seggpt-vit-large} & same checkpoint on both nodes \\
Backend service & \texttt{agriapp-local.service} & \texttt{agriapp-server.service} & separate lifecycle and logs \\
Android app package & \texttt{it.unipg.agriapp} & n/a & client package namespace \\
\bottomrule
\end{longtable}
